% =====================================================================
% §4 Diagnosis-Preserving Enhancement (C1 — hardest main leg)
% Nine-section reorder, W3 draft. Migrated/rewritten from main.tex L189-224.
% Serves Claim C1; ACCEPTANCE C1.1-C1.5. caveman OFF.
% Anti-drift locks per STORY: deterministic restoration not generative; FiLM positioned for
%   diagnostic fidelity not pixel quality; completeness partially recoverable per R1;
%   Prop3 carried by E7 plus non-vacuity condition; derive/sketch wording per R2.
% Macros assumed defined in preamble: \visienhance \Tw \qbar \Ldp \KL \Zenh \Zref
%   \Hent \tauenh \tauhigh \E ; appendix labels app:a2 / app:a2-prop3.
% =====================================================================

\section{Diagnosis-Preserving Enhancement}\label{sec:enhance}

\paragraph{Degradation taxonomy.}
Of the five quality dimensions, four (sharpness, brightness, colour temperature, and contrast)
admit deterministic restoration, whereas completeness ($q_3$, a truncated field of view) is only
\emph{partially recoverable} --- enhancement cannot synthesise tissue that was never captured, and we
make no claim that it does so. The agent therefore treats low completeness as a query-for-retake
signal rather than an enhancement target. This split --- which dimensions are recoverable and
\emph{within what severity window} --- is the design boundary that \cref{sec:enhance} formalises and
\cref{sec:exp} measures.

\subsection{Architecture: NAFNet with quality-conditional FiLM}\label{sec:arch}
\visienhance{} is a NAFNet backbone \citep{chen2022simple} augmented with FiLM modulation
\citep{perez2018film} conditioned on the per-input quality vector, so that restoration strength adapts
to \emph{how} an image is degraded rather than applying a fixed correction. The restoration map
$\Tw(x,q)$ is \emph{deterministic}: we deliberately avoid generative restoration (diffusion or
adversarial backbones), because a generative model can hallucinate diagnostic structure --- inventing a
pigment network or vascular pattern that was never present --- which is unacceptable in a clinical
setting (anti-drift rules~8--9; \cref{sec:discussion} analyses why diffusion fails here). We stress that
FiLM is \emph{not} a fidelity mechanism: its empirical role, established in \cref{sec:exp}~(E8), is
quality-conditional \emph{diagnostic} fidelity, not perceptual sharpening, and removing it leaves
per-image PSNR essentially unchanged.

\subsection{DP-Loss: a feature-level diagnosis-preserving objective}\label{sec:dploss}
Training minimises
\begin{equation}
\mathcal{L}=\mathcal{L}_1
  +\lambda_{\mathrm{LPIPS}}\,\mathcal{L}_{\mathrm{LPIPS}}
  +\lambda_{\mathrm{DP}}\,\Ldp
  +\lambda_{\mathrm{q}}\,\mathcal{L}_{\mathrm{hinge}},
\label{eq:total-loss}
\end{equation}
where $\mathcal{L}_1$ and $\mathcal{L}_{\mathrm{LPIPS}}$ govern pixel and perceptual fidelity, and the
diagnosis-preserving term is a \emph{feature-level} KL,
\begin{equation}
\Ldp \;=\; \KL\!\big(p_\phi^{\mathrm{enh}} \,\big\|\, p_\phi^{\mathrm{ref}}\big),
\label{eq:dploss}
\end{equation}
computed in the latent space of a frozen diagnostic feature extractor (the $1536\times7\times7$
penultimate features of an independent EfficientNet-B3 probe). Crucially, $\Ldp$ aligns the enhanced
image with its clean reference \emph{in the space that drives the diagnosis}, not in pixel space: a
restorer can be pixel-faithful yet still move the diagnostic posterior, and \eqref{eq:dploss} is what
penalises exactly that movement. The quality hinge $\mathcal{L}_{\mathrm{hinge}}$ encourages
$\qbar(\Tw(x,q))>\qbar(x)$ within the salvage band, tying restoration to a measured quality gain rather
than an assumed one.

\subsection{Mutual-information preservation (Lemma~\ref{lem:dp}) and a directional entropy result
(Proposition~\ref{prop:entropy})}\label{sec:prop3lemma3}
% Compact statements; full proofs in Appendix A2. R2: derive/sketch wording, not "prove".
We connect the DP objective \eqref{eq:dploss} to diagnosis through two results, stated compactly here
and derived in full in Appendix~\ref{app:a2}. The central guarantee is an information lower bound: under
a KL budget $\Ldp\leq\epsilon$, the mutual information between the enhanced latent and the label stays
within $\beta\sqrt{\epsilon}$ of the reference, $I(\Zenh;Y)\geq I(\Zref;Y)-\beta\sqrt{\epsilon}$, with
$\beta=M L_{q_\theta}/\sqrt{2}\approx0.74$ in the binary case ($M=\log K$, $L_{q_\theta}$ the
total-variation Lipschitz constant of the probe); the $\sqrt{\epsilon}$ rate is Pinsker-optimal for a KL
budget (Lemma~\ref{lem:dp}). In words: \emph{driving the feature-level KL small provably caps how much
diagnostic information enhancement can destroy.} The empirical counterpart of this bound is experiment
E7 (\cref{sec:exp}), which is the evidence we rely on.

We also derive a directional companion result on predictive entropy (Proposition~\ref{prop:entropy}):
the variance-driven component of the expected entropy carries a negative sign under quality improvement.
We are explicit that this is a \emph{directional characterisation of one term}, not a net claim --- a
quality-conditioned probe finds the net predictive entropy does \emph{not} fall after enhancement, so we
do \emph{not} present entropy reduction as an empirical result. The information-preservation role in the
paper is carried instead by Lemma~\ref{lem:dp} (validated by E7) together with the Stage-1
PSNR$\,\geq30$\,dB non-vacuity condition that keeps the bound non-empty
(Appendix~\ref{app:a2-prop3}). We state both results as derivations under listed assumptions
(``we derive'', ``under assumption'') rather than as watertight proofs.

% Compact statements of Prop 3 + Lemma 3 for the main body (§4.4).
% Full proofs: Appendix A2.1 / A2.2 (appendix/A2_prop3_lemma3.tex).

\begin{proposition}[Enhancement Entropy Reduction]\label{prop:entropy}
Under Assumptions A1--A4 (quality improvement, calibrated scoring, monotone prior,
diagnostic preservation), for $x$ in the salvage band and $q'=f_\eta(\Tw(x,q))$,
\begin{equation}
\E_{x}\!\big[\Hent(\hat{p}_{\Tw(x,q)})\big] \;\leq\;
\E_{x}\!\big[\Hent(\hat{p}_x)\big]
- \tfrac{d}{2}\log\tfrac{\sigma^2(\qbar)}{\sigma^2(\qbar')} + 2\beta\sqrt{\epsilon},
\label{eq:prop3}
\end{equation}
where the first negative term is the \qvib{} entropy gap (Lemma~1, strictly positive
since $\qbar'>\qbar \Rightarrow \sigma^2$ shrinks) and the residual $2\beta\sqrt{\epsilon}$
is the diagnostic-preservation slack (Lemma~\ref{lem:dp}).
\end{proposition}

\begin{lemma}[Diagnosis-Preserving Mutual-Information Bound]\label{lem:dp}
Let $q_\theta$ be $\Lq$-Lipschitz in total variation and let the label entropy be
bounded by $M=\log K$. If $\Ldp := \KL(p_\phi^{\mathrm{enh}} \,\|\, p_\phi^{\mathrm{ref}}) \leq \epsilon$, then
\begin{equation}
I(\Zenh; Y) \;\geq\; I(\Zref; Y) - \beta\sqrt{\epsilon},
\qquad \beta = \frac{M\,\Lq}{\sqrt{2}}.
\label{eq:lemma3}
\end{equation}
The $\sqrt{\epsilon}$ rate is Pinsker-optimal; a linear-$\epsilon$ bound would require a
stronger $\chi^2$-divergence assumption (Appendix~A2.2.5). The proof combines Pinsker's
inequality with a Lipschitz coupling argument (Appendix~A2.2).
\end{lemma}


\subsection{Three-stage training protocol}\label{sec:training}
We train $\visienhance$ in three stages, each isolating one of the terms in \eqref{eq:total-loss}:
\textbf{(1)} pixel restoration ($\mathcal{L}_1$ + LPIPS) trained to a PSNR gate, which also guards the
non-vacuity condition of Lemma~\ref{lem:dp} / Proposition~\ref{prop:entropy} (a restorer that fails the
PSNR floor cannot be claimed to preserve information; Appendix~\ref{app:a2-prop3});
\textbf{(2)} DP-Loss fine-tuning that tightens the feature-level KL $\Ldp$ toward the budget $\epsilon$,
the stage whose ablation is E7; and
\textbf{(3)} a quality-hinge stage that calibrates the $\qbar$ improvement inside the salvage band
$[\tauenh,\tauhigh]$, so that enhancement is only credited where it measurably raises quality.
Stage~(2) is the stage that turns the Lemma~\ref{lem:dp} budget into an optimisation target; the
empirical effect of switching it on versus off is reported as the DP-Loss ablation E7 in
\cref{sec:exp}.
